
tères pour que \(F(k)\longrightarrow F(K)\) soit injectif (resp. surjectif).

\textbf{7.18.2.} Supposons que l'on soit dans l'un des deux cas suivants :
\begin{enumerate}[label=\alph*)]
\item Le corps \(k\) est séparablement clos, et \(K/k\) est une extension séparable.
\item Le corps \(k\) est infini, et \(K\) en est une extension pure.
\end{enumerate}
Alors l'application \(F(k)\to F(K)\) est injective.

En effet, utilisant le fait que \(K\) est limite inductive filtrante de sous-algèbres de type fini sur \(k\), avec des morphismes de transition plats, et la commutation de \(F\) aux limites inductives filtrantes, on est ramené dans le cas a) à prouver que si \(A\) est une algèbre intègre de type fini sur \(k\), dont le corps des fractions \(K\) est séparable sur \(k\), alors \(F(k)\to F(A)\) est injectif. Or les hypothèses faites sur \(k\) et sur \(K\) impliquent qu'il existe un \(k\)-homomorphisme \(A\to k\); alors le composé \(F(k)\to F(A)\to F(k)\) est l'identité, d'où la conclusion. Dans le cas b), on se ramène de même, par passage à la limite, à prouver que si \(A\) est une algèbre intègre de type fini sur \(k\) telle que son corps des fractions soit une extension pure de \(k\), alors \(F(k)\to F(A)\) est injectif. Comme \(k\) est infini, l'hypothèse implique encore qu'il existe un \(k\)-homomorphisme \(A\to k\), et on achève comme ci-dessus.

\textbf{7.18.3.} Supposons maintenant que \(F\) soit un foncteur à valeurs dans la catégorie \((\mathrm{Ab})\) des groupes abéliens, et supposons que pour toute extension finie \(k'\) de \(k\), on ait défini un homomorphisme norme
\[
  N_{k'/k}:F(k')\longrightarrow F(k)
\]
tel que le composé
\[
  F(k)\longrightarrow F(k')\xrightarrow{N_{k'/k}}F(k)
\]
soit la multiplication par \([k':k]\). Sous ces conditions, on a :
\begin{enumerate}[label=(\roman*)]
